% Book-style model test: page 1 = MCQ (3 columns), page 2 = CQ (2 columns)
\documentclass[10pt]{article}
\usepackage[tiro]{edtech}
\geometry{margin=11mm,top=10mm,bottom=12mm}
\etset{brand=EduBrand, logo=logo-placeholder.png}
\pagestyle{empty}
\setlength{\columnseprule}{0.4pt}
\setlength{\columnsep}{5mm}
\begin{document}
%% ------------------------------------------------------------ MCQ page
\noindent\etlogo[0.8cm]\hfill{\small[সেট : ক]}\par\vspace{-2pt}
\modeltestheader{এক্সক্লুসিভ মডেল টেস্ট ০১}{রসায়ন প্রথম পত্র\etsep বহুনির্বাচনি অভীক্ষা}{২৫ মিনিট}{২৫}{176}
 {সরবরাহকৃত বহুনির্বাচনি অভীক্ষার উত্তরপত্রে প্রশ্নের ক্রমিক নম্বরের বিপরীতে প্রদত্ত বর্ণসংবলিত বৃত্তসমূহ হতে সঠিক/সর্বোৎকৃষ্ট উত্তরের বৃত্তটি বল পয়েন্ট কলম দ্বারা সম্পূর্ণ ভরাট কর। প্রতিটি প্রশ্নের মান ১।}
\optstyle{circled}\qnumstyle{dot}
\begin{multicols}{3}\small
\mcq[খ]{কোনটি ক্ষারীয় বাফার?}{\ce{CH3COOH + CH3COONa}}{\ce{NH4OH + NH4Cl}}{\ce{H2CO3 + NaHCO3}}{\ce{NaH2PO4 + Na2HPO4}}
\mcq[গ]{$3d$ অরবিটালের জন্য $l$ এর মান কত?}{0}{1}{2}{3}
\mcq[ঘ]{কোনটির প্রথম আয়নিকরণ শক্তি সবচেয়ে বেশি?}{Li}{Be}{B}{N}
\mcq[ক]{pH = 3 দ্রবণে $[\ce{H+}]$ কত?}{$10^{-3}$ M}{$10^{-11}$ M}{3 M}{$10^{3}$ M}
\mcqmulti[ক]{$sp^3$ সংকরণ ঘটে—}{\ce{CH4} অণুতে}{\ce{NH3} অণুতে}{\ce{BF3} অণুতে}
\mcq[খ]{কোনটি জারণ-বিজারণ বিক্রিয়া নয়?}{\ce{Zn + CuSO4 -> ZnSO4 + Cu}}{\ce{HCl + NaOH -> NaCl + H2O}}{\ce{2Mg + O2 -> 2MgO}}{\ce{H2 + Cl2 -> 2HCl}}
\mcq[খ]{হাইড্রোজেন বর্ণালির বামার সিরিজ কোন অঞ্চলে পড়ে?}{অতিবেগুনি}{দৃশ্যমান}{অবলোহিত}{দূর অবলোহিত}
\mcq[খ]{\ce{N2 + 3H2 <=> 2NH3} বিক্রিয়ায় $\Delta n$ এর মান কত?}{2}{$-2$}{0}{1}
\mcq[ক]{সর্বোচ্চ তড়িৎ ঋণাত্মক মৌল কোনটি?}{F}{O}{Cl}{N}
\mcq[খ]{কোনটি উভধর্মী অক্সাইড?}{\ce{Na2O}}{\ce{Al2O3}}{\ce{SO3}}{\ce{CaO}}
\stimulus{A, B ও C মৌলের পারমাণবিক সংখ্যা যথাক্রমে 11, 17 ও 20।}
\mcq[ক]{A ও B দ্বারা গঠিত যৌগে কোন বন্ধন বিদ্যমান?}{আয়নিক}{সমযোজী}{ধাতব}{সন্নিবেশ}
\mcq[খ]{C মৌলটি পর্যায় সারণির কোন গ্রুপে অবস্থিত?}{1}{2}{13}{17}
\mcq[গ]{\SI{25}{\degreeCelsius} তাপমাত্রায় পানির আয়নিক গুণফল কত?}{$10^{-7}$}{$10^{-12}$}{$10^{-14}$}{$10^{-15}$}
\mcq[গ]{কোনটি প্রাইমারি স্ট্যান্ডার্ড পদার্থ?}{\ce{NaOH}}{\ce{HCl}}{\ce{Na2CO3}}{\ce{KMnO4}}
\mcqmulti[ক]{এসিড বৃষ্টির জন্য দায়ী—}{\ce{SO2}}{\ce{NO2}}{\ce{CH4}}
\mcq[ক]{\ce{CO2} অণুর আকৃতি কেমন?}{সরলরৈখিক}{কৌণিক}{ত্রিকোণাকার}{চতুস্তলকীয়}
\mcq[খ]{$n=3$ শক্তিস্তরে সর্বোচ্চ কতটি ইলেকট্রন থাকতে পারে?}{8}{18}{32}{9}
\mcq[খ]{\ce{NH3} অণুতে নিঃসঙ্গ ইলেকট্রন যুগল কয়টি?}{0}{1}{2}{3}
\mcq[ক]{দ্রাব্যতা গুণফলের একক $\mathrm{mol^2\,L^{-2}}$ কোনটির?}{\ce{AgCl}}{\ce{PbI2}}{\ce{Ca3(PO4)2}}{\ce{Fe(OH)3}}
\mcq[ক]{\ce{N2 + 3H2 <=> 2NH3}; চাপ বৃদ্ধি করলে সাম্যাবস্থা কোন দিকে সরবে?}{ডানে}{বামে}{অপরিবর্তিত}{প্রথমে বামে পরে ডানে}
\mcq[খ]{\SI{0.1}{M} \ce{HCl} দ্রবণের pH কত?}{0.1}{1}{13}{2}
\mcq[খ]{কোনটি লুইস এসিড?}{\ce{NH3}}{\ce{BF3}}{\ce{H2O}}{\ce{OH-}}
\mcq[খ]{শিখা পরীক্ষায় সোডিয়াম কোন বর্ণ দেয়?}{ইটলাল}{সোনালি হলুদ}{হালকা বেগুনি}{আপেল সবুজ}
\mcq[খ]{তীব্র এসিড ও মৃদু ক্ষারের টাইট্রেশনে উপযুক্ত নির্দেশক কোনটি?}{ফেনলফথ্যালিন}{মিথাইল অরেঞ্জ}{লিটমাস}{কোনোটিই নয়}
\mcq[ঘ]{কোনটি অপোলার অণু?}{\ce{H2O}}{\ce{NH3}}{\ce{HCl}}{\ce{CCl4}}
\end{multicols}
\vspace{-4pt}\answerkey

%% ------------------------------------------------------------ CQ page
\newpage
\setcounter{etq}{0}\cqstyle{book}
\modeltestheader{এক্সক্লুসিভ মডেল টেস্ট ০১}{রসায়ন প্রথম পত্র\etsep সৃজনশীল প্রশ্ন}{২ ঘণ্টা ৩৫ মিনিট}{৫০}{176}
 {ডান পাশের সংখ্যা প্রশ্নের পূর্ণমান জ্ঞাপক। প্রদত্ত উদ্দীপকগুলো মনোযোগ সহকারে পড় এবং সংশ্লিষ্ট প্রশ্নগুলোর উত্তর দাও। যে কোনো \textbf{পাঁচটি} প্রশ্নের উত্তর দিতে হবে।}
\begin{multicols}{2}\small
\cq{হাইড্রোজেন পরমাণুর একটি ইলেকট্রন চিত্রানুযায়ী শক্তিস্তরে স্থানান্তরিত হয়:\par
\begin{center}\begin{tikzpicture}[xscale=0.9,yscale=0.95,font=\scriptsize]
 \foreach \n/\y in {1/0,2/1.5,3/2.2,4/2.6}{\draw[thick](0,\y)--(3.2,\y) node[right]{$n=\n$};}
 \draw[-{Latex[length=2mm]},thick](0.8,2.6)--(0.8,1.5) node[midway,left]{X};
 \draw[-{Latex[length=2mm]},thick](2.2,2.2)--(2.2,0) node[midway,right]{Y};
\end{tikzpicture}\end{center}
[$R_H=\SI{1.097e7}{m^{-1}}$]}
{অরবিটাল কাকে বলে?}
{$2d$ অরবিটাল সম্ভব নয় কেন? ব্যাখ্যা কর।}
{X স্থানান্তরে বিকিরিত রশ্মির তরঙ্গদৈর্ঘ্য নির্ণয় কর।}
{X ও Y স্থানান্তরের মধ্যে কোনটি খালি চোখে দেখা যাবে? গাণিতিকভাবে বিশ্লেষণ কর।}

\cq{\begin{center}\renewcommand\arraystretch{1.2}\begin{tabular}{|c|c|}\hline
 মৌল & শেষ শক্তিস্তরের ইলেকট্রন বিন্যাস\\\hline
 A & $ns^2np^3$\\\hline B & $(n+1)s^2(n+1)p^3$\\\hline C & $(n+1)s^2(n+1)p^5$\\\hline
\end{tabular}\quad এখানে $n=2$\end{center}}
{হুন্ডের নীতি কী?}
{\ce{AlCl3} সমযোজী হলেও \ce{AlF3} আয়নিক কেন?}
{\ce{AC5} ও \ce{BC5} গঠন সম্ভব কি না ব্যাখ্যা কর।}
{A ও B এর হাইড্রাইডের সংকরণ একই হলেও বন্ধন কোণ ভিন্ন হওয়ার কারণ বিশ্লেষণ কর।}

\cq{নিচের উদ্দীপকটি লক্ষ কর:\par
\begin{center}\begin{tikzpicture}[font=\scriptsize]
 \node[draw,thick,minimum width=2.3cm,minimum height=1cm,align=center](a){\SI{0.05}{M} \ce{H2SO4}\\দ্রবণ};
 \node[draw,thick,minimum width=3.2cm,minimum height=1cm,align=center,right=4mm of a](b){\ce{CH3COOH} ও \ce{CH3COONa}\\এর মিশ্রিত দ্রবণ};
 \node[below=1mm of a]{A-পাত্র};\node[below=1mm of b]{B-পাত্র};
\end{tikzpicture}\end{center}}
{আয়নিকরণ বিভব কী?}
{\ce{HNO3} ও \ce{H3PO4} এর মধ্যে কোনটি তীব্র এসিড? ব্যাখ্যা কর।}
{উদ্দীপকের A পাত্রের দ্রবণের pOH হিসাব কর।}
{B পাত্রের দ্রবণে সামান্য তীব্র এসিড বা ক্ষার যোগ করলে pH পরিবর্তিত হবে কি না বিশ্লেষণ কর।}

\cq{\ce{N2(g) + 3H2(g) <=> 2NH3(g)};\ $\Delta H=\SI{-92}{kJ.mol^{-1}}$। \SI{500}{\degreeCelsius} তাপমাত্রায় $K_c=\SI{6.0e-2}{mol^{-2}.L^{2}}$।}
{সাম্যধ্রুবক কাকে বলে?}
{তাপমাত্রা বাড়ালে উদ্দীপকের বিক্রিয়ার উৎপাদ কমে কেন?}
{উক্ত তাপমাত্রায় $K_p$ এর মান নির্ণয় কর।}
{হেবার পদ্ধতিতে উদ্দীপকের বিক্রিয়ার সর্বোত্তম উৎপাদনের শর্তগুলো লা-শাতেলিয়ের নীতির আলোকে বিশ্লেষণ কর।}

\cq{\begin{center}\begin{tikzpicture}[font=\scriptsize]
 \draw[thick](0,1.4)--(0,0)--(1.8,0)--(1.8,1.4); \node[align=center] at (0.9,0.65){\SI{30}{mL}\\\SI{5e-3}{M}\\\ce{YCl2}};
 \draw[thick](3,1.4)--(3,0)--(4.8,0)--(4.8,1.4); \node[align=center] at (3.9,0.65){\SI{20}{mL}\\\SI{3e-3}{M}\\\ce{NaF}};
 \node at (2.4,0.7){\SI{25}{\degreeCelsius}};
 \node[below] at (0.9,0){১ নং দ্রবণ}; \node[below] at (3.9,0){২ নং দ্রবণ};
\end{tikzpicture}\end{center}
[Y এর পারমাণবিক ভর 40]}
{হাইড্রোজেন বন্ধন কাকে বলে?}
{পানির ডাইপোল মোমেন্ট থাকার কারণ ব্যাখ্যা কর।}
{১নং দ্রবণের ক্যাটায়নের শনাক্তকরণ পরীক্ষা সমীকরণসহ লেখ।}
{১নং ও ২নং দ্রবণ একত্রে মিশ্রিত করলে \ce{YF2} এর অধঃক্ষেপ পড়বে কি না গাণিতিকভাবে বিশ্লেষণ কর। [\ce{YF2} এর $K_{sp}=4\times10^{-11}$]}

\cq{খেজুরের রস $\longrightarrow$ A (খাদ্য সংরক্ষক)। A একটি সোডিয়াম লবণ যা কোমল পানীয়তে ব্যবহৃত হয়।}
{পোলারায়ন কাকে বলে?}
{সোডিয়ামের দ্বিতীয় আয়নিকরণ শক্তি অনেক বেশি কেন?}
{খাদ্য সংরক্ষক A এর প্রস্তুতি সমীকরণসহ লেখ।}
{A যৌগটির খাদ্য সংরক্ষক হিসেবে উপযুক্ততা বিশ্লেষণ কর।}

\cq{একজন শিক্ষার্থী \SI{250}{mL} আয়তনিক ফ্লাস্কে \SI{1.325}{g} অনার্দ্র \ce{Na2CO3} দ্রবীভূত করে একটি প্রমাণ দ্রবণ তৈরি করল।}
{মোলারিটি কাকে বলে?}
{প্রমাণ দ্রবণ তৈরিতে আয়তনিক ফ্লাস্ক ব্যবহার করা হয় কেন?}
{দ্রবণটির মোলারিটি নির্ণয় কর।}
{দ্রবণটির \SI{20}{mL} সম্পূর্ণ প্রশমিত করতে \SI{0.1}{M} \ce{HCl} এর কত আয়তন প্রয়োজন— বিশ্লেষণ কর।}

\cq{X ও Y দুটি মৌলের পারমাণবিক সংখ্যা যথাক্রমে 24 ও 29।}
{আউফবাউ নীতি কী?}
{\ce{Cu^2+} আয়ন রঙিন কিন্তু \ce{Zn^2+} আয়ন বর্ণহীন কেন?}
{X এর ইলেকট্রন বিন্যাস ব্যতিক্রমী কেন— ব্যাখ্যা কর।}
{X ও Y এর মধ্যে কোনটির জটিল যৌগ গঠনের প্রবণতা বেশি? যুক্তিসহ বিশ্লেষণ কর।}
\end{multicols}
\end{document}
